\chapter{市场分析}

\section{目标市场定位}

本项目的目标市场为\hlgreen{中国小农户病虫害诊断与防治的数字化服务市场}，具体聚焦：以水稻、小麦、玉米等主粮作物及高价值果蔬为重点作物，以规模在 50 亩以下的小农户与适度规模经营主体为核心用户，以县、镇两级为基本渗透单元，向全国主要农业产区逐步辐射。

\section{市场规模与增长潜力}

据公开数据，我国智慧农业市场规模 2024 年已达约 8500 亿元，预计 2028 年将突破 1.4 万亿元，年均复合增速超过 13\%。其中，植保环节作为农业生产中「投入高频、损失敏感、知识门槛高」的核心环节，是数字化改造价值最高的细分场景之一。即便仅按智慧农业市场 5\% 的植保数字化渗透率估算，病虫害诊断防治相关服务的潜在市场规模也达数百亿元量级。

\begin{figure}[H]
  \centering
  \includegraphics[width=0.86\linewidth]{market_size.png}
  \caption{中国智慧农业市场规模及增速（亿元，预测）}
\end{figure}

\section{目标用户画像与需求特征}

我国有超过 2 亿户小农户，其共同特征为：\hlgreen{规模小、分散、老龄化程度高、专业植保知识薄弱、价格敏感但对减产风险极度敏感}。具体而言：

\begin{itemize}[leftmargin=1.6em]
  \item \hlgreen{作物结构}：以水稻、小麦、玉米三大主粮为主，果蔬等高价值经济作物占比提升。
  \item \hlgreen{行为习惯}：智能手机普及率已达 76\%，但农事决策仍高度依赖熟人经验与农资店推荐。
  \item \hlgreen{核心诉求}：要的不是「复杂的农业百科全书」，而是「看到病害照片就能立刻告诉我怎么办」的确定性答案。
  \item \hlgreen{付费意愿}：对「能直接挽回损失」的精准防治方案付费意愿强，对「锦上添花」的泛内容付费意愿弱。
\end{itemize}

\begin{figure}[H]
  \centering
  \includegraphics[width=0.86\linewidth]{user_profile.png}
  \caption{目标用户作物结构}
\end{figure}

\section{竞争格局分析}

\subsection{主流竞品类型梳理}

当前市场参与者可分为三类：一是\hlgreen{传统农技服务体系}（基层农技站、农资店「土专家」），线下为主、经验驱动；二是\hlgreen{综合型农业 APP}（农资电商、农业资讯、农技问答类平台），功能大而全但门槛高；三是\hlgreen{少数植保 AI 试点产品}，多停留在实验室或政府示范阶段，尚未形成规模化、面向小农户的成熟产品。

\subsection{核心竞品优劣势对比}

\begin{figure}[H]
  \centering
  \includegraphics[width=0.80\linewidth]{competitive_radar.png}
  \caption{竞争格局雷达对比}
\end{figure}

\subsection{本项目的差异化竞争优势}

\begin{enumerate}[leftmargin=1.6em]
  \item \hlgreen{极致轻量}：单一功能闭环，端侧离线可用，零学习成本，直接面向小农户，而非「面向农业从业者的专业工具」。
  \item \hlgreen{方案落地}：不止识别，更输出「用什么药、怎么治、什么时候防」的可执行方案，价值直达生产决策。
  \item \hlgreen{本地化预警}：基于 LBS 与当地气象、历史发病数据，主动推送当季高发病虫害预防提示，变「被动求医」为「主动防御」。
  \item \hlgreen{渠道壁垒}：借力农资终端、村镇基层组织与基层农技体系，构建竞争对手难以复制的线下触达网络。
\end{enumerate}

\section{市场空白点与机会总结}

综上，市场存在一个清晰的\hlgold{「三不管」空白带}：传统农技服务「管不过来」、综合农业 APP「不愿做」、植保 AI 试点「没落地」。小农户「拍照识病、即时拿方案、提前防病」这一刚需，至今没有一个真正好用、用得起、用得惯的产品来满足。率先以极简产品切入、以线下渠道起量、以数据壁垒沉淀，本项目有望在 3 年内成为「病虫害 AI 诊断」这一垂直品类的用户心智第一选择。
